\section*{Примечания о чтениях}
\subsection*{\texttt{GELFAND-P001455}}
\par\noindent\leavevmode\phantomsection\label{GELFAND-0158-U145501}\label{GELFAND-P001455:configuration-cancel}\textbf{\texttt{GELFAND-0158-U145501}} (configuration-cancel).\nopagebreak[4]\par
Между «Фуксу» и «не» один знак густо зачёркнут; видны верхний наклонный штрих и нижний возврат.\par
Кандидаты: у; д. Чтение не установлено в дипломатическом исходнике.\par
Исходные координаты [1310,827,1520,930], увеличение 7x; также строка [940,730,2030,960], 3x. Зачёркнутое начало не восстановлено по смыслу.\par
Повторная сверка: Повторно просмотрены исходник и увеличения 3×, 5× и 10×. Сохранены исходные наблюдения, варианты чтения и неуверенности; математическое рассуждение не заменяет буквальное чтение. Частичные воссоздания рисунков не утверждают неустановленных соединений или индексов.\par\medskip
\par\noindent\leavevmode\phantomsection\label{GELFAND-0158-U145502}\label{GELFAND-P001455:task-number}\textbf{\texttt{GELFAND-0158-U145502}} (task-number).\nopagebreak[4]\par
После подчёркнутого «Задача» знак с петлёй и длинным верхним штрихом, затем точка.\par
Кандидаты: 6; в. Чтение не установлено в дипломатическом исходнике.\par
Исходные координаты [480,1100,600,1210], 6x; строка [200,1100,2150,1510], 3x. Вопрос о цифре или букве оставлен открытым.\par
Повторная сверка: Повторно просмотрены исходник и увеличения 3×, 5× и 10×. Сохранены исходные наблюдения, варианты чтения и неуверенности; математическое рассуждение не заменяет буквальное чтение. Частичные воссоздания рисунков не утверждают неустановленных соединений или индексов.\par\medskip
\par\noindent\leavevmode\phantomsection\label{GELFAND-0158-U145503}\label{GELFAND-P001455:vertex-count-symbol}\textbf{\texttt{GELFAND-0158-U145503}} (vertex-count-symbol).\nopagebreak[4]\par
В определении числа вершин и формулах (13)–(16′) повторяется знак, сходный с рукописными ν и v, с индексом i или числом.\par
Кандидаты: \textbackslash{}nu\_i; v\_i. Чтение не установлено в дипломатическом исходнике.\par
Проверены [270,1515,2160,2190], 3x, и [115,2200,1730,3150], 3x. В следующих формулах использован представитель ν; выбор алфавита остаётся неуверенным. Отдельный знак в определении дополнительно проверен на [765,1935,910,2065], 10x; различие ν/v не снято.\par
Повторная сверка: Повторно просмотрены исходник и увеличения 3×, 5× и 10×. Сохранены исходные наблюдения, варианты чтения и неуверенности; математическое рассуждение не заменяет буквальное чтение. Частичные воссоздания рисунков не утверждают неустановленных соединений или индексов.\par\medskip
\par\noindent\leavevmode\phantomsection\label{GELFAND-0158-U145504}\label{GELFAND-P001455:cancelled-equation-number}\textbf{\texttt{GELFAND-0158-U145504}} (cancelled-equation-number).\nopagebreak[4]\par
Слева от последней формулы есть густо зачёркнутый номер в скобках; ниже/правее читается (16′).\par
Кандидаты: (17); (16). Чтение не установлено в дипломатическом исходнике.\par
Исходные координаты [120,3290,355,3410], 7x; также [100,3190,1920,3440], 3x. Читаемый номер (16′) сохранён независимо от зачёркнутого.\par
Повторная сверка: Повторно просмотрены исходник и увеличения 3×, 5× и 10×. Сохранены исходные наблюдения, варианты чтения и неуверенности; математическое рассуждение не заменяет буквальное чтение. Частичные воссоздания рисунков не утверждают неустановленных соединений или индексов.\par\medskip
\subsection*{\texttt{GELFAND-P001456}}
\par\noindent\leavevmode\phantomsection\label{GELFAND-0158-U145601}\label{GELFAND-P001456:initial-charge-cancel}\textbf{\texttt{GELFAND-0158-U145601}} (initial-charge-cancel).\nopagebreak[4]\par
Перед 60(6−i) зачёркнуто короткое выражение; i в конце различим, первая часть перекрыта несколькими горизонтальными штрихами.\par
Кандидаты: 60-i; 6-i. Чтение не установлено в дипломатическом исходнике.\par
Исходные координаты [870,1280,1110,1375], 8x, и [75,1045,2230,1555], 3x. Читаемая замена 60(6−i) сохранена отдельно.\par
Повторная сверка: Повторно просмотрены исходник и увеличения 3×, 5× и 10×. Сохранены исходные наблюдения, варианты чтения и неуверенности; математическое рассуждение не заменяет буквальное чтение. Частичные воссоздания рисунков не утверждают неустановленных соединений или индексов.\par\medskip
\par\noindent\leavevmode\phantomsection\label{GELFAND-0158-U145602}\label{GELFAND-P001456:impossible-cancel}\textbf{\texttt{GELFAND-0158-U145602}} (impossible-cancel).\nopagebreak[4]\par
После «невозможно» перед «ввиду» зачёркнут один небольшой знак или короткое начало слова с длинным нижним штрихом.\par
Кандидаты: не записаны. Чтение не установлено в дипломатическом исходнике.\par
Исходные координаты [1160,1785,1235,1900], 8x; контекст [65,1565,2210,2150], 3x. Под зачёркиванием надёжное чтение не установлено.\par
Повторная сверка: Повторно просмотрены исходник и увеличения 3×, 5× и 10×. Сохранены исходные наблюдения, варианты чтения и неуверенности; математическое рассуждение не заменяет буквальное чтение. Частичные воссоздания рисунков не утверждают неустановленных соединений или индексов.\par\medskip
\par\noindent\leavevmode\phantomsection\label{GELFAND-0158-U145603}\label{GELFAND-P001456:situation-one-cancel}\textbf{\texttt{GELFAND-0158-U145603}} (situation-one-cancel).\nopagebreak[4]\par
Между «0 или 1» и «верш.» есть зачёркнутое слово, начинающееся р и содержащее высокий штрих б.\par
Кандидаты: рёбр; рёбер. Чтение не установлено в дипломатическом исходнике.\par
Исходные координаты [1015,2315,1185,2425], 8x; строка [55,2170,2010,2580], 3x. Конец слова перекрыт штрихами.\par
Повторная сверка: Повторно просмотрены исходник и увеличения 3×, 5× и 10×. Сохранены исходные наблюдения, варианты чтения и неуверенности; математическое рассуждение не заменяет буквальное чтение. Частичные воссоздания рисунков не утверждают неустановленных соединений или индексов.\par\medskip
\par\noindent\leavevmode\phantomsection\label{GELFAND-0158-U145604}\label{GELFAND-P001456:cancelled-enumeration}\textbf{\texttt{GELFAND-0158-U145604}} (cancelled-enumeration).\nopagebreak[4]\par
Между пояснением к (23) и строкой (24) отдельно зачёркнут знак с петлеобразным верхом и нижним пересечением.\par
Кандидаты: 2. Чтение не установлено в дипломатическом исходнике.\par
Исходные координаты [330,2580,410,2685], 8x. Возможная начатая цифра 2; её назначение не утверждается.\par
Повторная сверка: Повторно просмотрены исходник и увеличения 3×, 5× и 10×. Сохранены исходные наблюдения, варианты чтения и неуверенности; математическое рассуждение не заменяет буквальное чтение. Частичные воссоздания рисунков не утверждают неустановленных соединений или индексов.\par\medskip
\par\noindent\leavevmode\phantomsection\label{GELFAND-0158-U145605}\label{GELFAND-P001456:situation-two-cancel}\textbf{\texttt{GELFAND-0158-U145605}} (situation-two-cancel).\nopagebreak[4]\par
После «соед.» зачёркнута группа, начинающаяся «ребр…» и оканчивающаяся отдельным с.\par
Кандидаты: ребра с; ребро с. Чтение не установлено в дипломатическом исходнике.\par
Исходные координаты [765,2700,1040,2820], 7x, и [65,2580,2070,2870], 3x. Читаемая последующая фраза «с ≤2 вершин кратности 5» оставлена как в источнике.\par
Повторная сверка: Повторно просмотрены исходник и увеличения 3×, 5× и 10×. Сохранены исходные наблюдения, варианты чтения и неуверенности; математическое рассуждение не заменяет буквальное чтение. Частичные воссоздания рисунков не утверждают неустановленных соединений или индексов.\par\medskip
\par\noindent\leavevmode\phantomsection\label{GELFAND-0158-U145606}\label{GELFAND-P001456:vertex-count-symbol}\textbf{\texttt{GELFAND-0158-U145606}} (vertex-count-symbol).\nopagebreak[4]\par
В (17) и (19) повторяется рукописный знак для числа вершин, допускающий ν и v.\par
Кандидаты: \textbackslash{}nu; v. Чтение не установлено в дипломатическом исходнике.\par
Верх страницы [70,30,1990,295], 3x, сопоставлен с определением на P001455. Здесь и далее в формуле выбран представитель ν без снятия неуверенности в алфавите.\par
Повторная сверка: Повторно просмотрены исходник и увеличения 3×, 5× и 10×. Сохранены исходные наблюдения, варианты чтения и неуверенности; математическое рассуждение не заменяет буквальное чтение. Частичные воссоздания рисунков не утверждают неустановленных соединений или индексов.\par\medskip
\par\noindent\leavevmode\phantomsection\label{GELFAND-0158-U145607}\label{GELFAND-P001456:diagram-18-right}\textbf{\texttt{GELFAND-0158-U145607}} (diagram-18-right).\nopagebreak[4]\par
У правого пятиугольного рисунка в (18) верхнее правое ребро выделено густым клиновидным штрихом; возле его правого конца есть выступ. Остальной контур, центральный узел и пять радиальных отрезков различимы.\par
Кандидаты: Утолщённое/обведённое ребро с остаточным штрихом; Ребро с дополнительной стрелочной или иной отметкой. Чтение не установлено в дипломатическом исходнике.\par
Исходные координаты [80,300,1680,710], увеличение 3x; верхняя правая сторона приблизительно [990,310,1110,390]. Характер выделения не выведен из теории графов. SourceFigure сохраняет описание всей третьей фигуры и подпись. Дополнительно проверено [990,305,1115,405], 10x: густой клин и выступ видны, назначение остаётся неуверенным.\par
Повторная сверка: Повторно просмотрены исходник и увеличения 3×, 5× и 10×. Сохранены исходные наблюдения, варианты чтения и неуверенности; математическое рассуждение не заменяет буквальное чтение. Частичные воссоздания рисунков не утверждают неустановленных соединений или индексов.\par\medskip
\subsection*{\texttt{GELFAND-P001457}}
\par\noindent\leavevmode\phantomsection\label{GELFAND-0158-U145701}\label{GELFAND-P001457:vanishing-cycles-cancel}\textbf{\texttt{GELFAND-0158-U145701}} (vanishing-cycles-cancel).\nopagebreak[4]\par
В пояснении справа от формулы (2) перед читаемым «циклов» густо зачёркнуто слово; под горизонтальными штрихами различимы несколько нижних петель и длинный правый спуск.\par
Кандидаты: коциклов; когомол.. Чтение не установлено в дипломатическом исходнике.\par
Проверены исходные координаты [1715,2550,1970,2670], 8x, и более широкий участок [1700,2550,2040,2690], 6x; контекст [115,2380,2360,2840], 3x. Слово перекрыто несколькими штрихами, чтения остаются предположительными.\par
Повторная сверка: Повторно просмотрены исходник и увеличения 3×, 5× и 10×. Сохранены исходные наблюдения, варианты чтения и неуверенности; математическое рассуждение не заменяет буквальное чтение. Частичные воссоздания рисунков не утверждают неустановленных соединений или индексов.\par\medskip
\par\noindent\leavevmode\phantomsection\label{GELFAND-0158-U145702}\label{GELFAND-P001457:hypercohomology-word}\textbf{\texttt{GELFAND-0158-U145702}} (hypercohomology-word).\nopagebreak[4]\par
Под условием ε≠0 мало в скобках написано слово, начинающееся «гипер…» и заканчивающееся «…гомологии»; отдельное «ко» между р и го не различается надёжно.\par
Кандидаты: гипергомологии; гиперкогомологии. Чтение не установлено в дипломатическом исходнике.\par
Исходные координаты [430,2730,910,2845], увеличение 6x; полный контекст [115,2380,2360,2840], 3x. Независимый повторный просмотр также не снял сомнение; ожидаемый термин не подставлен по математическому смыслу.\par
Повторная сверка: Повторно просмотрены исходник и увеличения 3×, 5× и 10×. Сохранены исходные наблюдения, варианты чтения и неуверенности; математическое рассуждение не заменяет буквальное чтение. Частичные воссоздания рисунков не утверждают неустановленных соединений или индексов.\par\medskip
\subsection*{\texttt{GELFAND-P001458}}
\par\noindent\leavevmode\phantomsection\label{GELFAND-0158-U145801}\label{GELFAND-P001458:top-cancellation}\textbf{\texttt{GELFAND-0158-U145801}} (top-cancellation).\nopagebreak[4]\par
В первой строке между «того,» и «здесь» слово зачёркнуто горизонтальными штрихами. Над ним ясно вставлено «при чём». Исходный кадр 2408×3480; область (795,30,1130,150).\par
Кандидаты: почему; зачем. Чтение не установлено в дипломатическом исходнике.\par
Область рассмотрена при 6×; начальные и средние буквы перекрыты. Кандидаты сохраняют именно зачёркнутый текст, не подменяют вставку.\par
Повторная сверка: Повторно просмотрены исходник и увеличения 3×, 5× и 10×. Сохранены исходные наблюдения, варианты чтения и неуверенности; математическое рассуждение не заменяет буквальное чтение. Частичные воссоздания рисунков не утверждают неустановленных соединений или индексов.\par\medskip
\par\noindent\leavevmode\phantomsection\label{GELFAND-0158-U145802}\label{GELFAND-P001458:vanishing-cycle-sketch}\textbf{\texttt{GELFAND-0158-U145802}} (vanishing-cycle-sketch).\nopagebreak[4]\par
Слева под первой строкой верхние дуги сходятся слева, справа находится узкий заштрихованный овал; ниже горизонтальный овал с двумя точками и соединяющим их штрихом с левой головкой. От верхнего правого овала вниз идёт прерывистая линия. Исходная область (285,145,1190,440).\par
Кандидаты: не записаны. Чтение не установлено в дипломатическом исходнике.\par
Рассмотрено при 3×. Из-за пересечений и утолщений нельзя уверенно воспроизвести все направления и геометрию; SourceFigure сохраняет связь с оригинальным рисунком. Подпись имеет отдельную reading-запись.\par
Повторная сверка: Повторно просмотрены исходник и увеличения 3×, 5× и 10×. Сохранены исходные наблюдения, варианты чтения и неуверенности; математическое рассуждение не заменяет буквальное чтение. Частичные воссоздания рисунков не утверждают неустановленных соединений или индексов.\par\medskip
\par\noindent\leavevmode\phantomsection\label{GELFAND-0158-U145803}\label{GELFAND-P001458:cycle-caption}\textbf{\texttt{GELFAND-0158-U145803}} (cycle-caption).\nopagebreak[4]\par
Подпись справа от рисунка: «цикл стя[н/г]. в точку». Последняя буква сокращения неразличима между н и г. Область (688,205,1155,313).\par
Кандидаты: стян.; стяг.. Чтение не установлено в дипломатическом исходнике.\par
Рассмотрено при 6×; независимый повторный просмотр также оставил два варианта с предпочтением стян. Сокращение не развёрнуто.\par
Повторная сверка: Повторно просмотрены исходник и увеличения 3×, 5× и 10×. Сохранены исходные наблюдения, варианты чтения и неуверенности; математическое рассуждение не заменяет буквальное чтение. Частичные воссоздания рисунков не утверждают неустановленных соединений или индексов.\par\medskip
\par\noindent\leavevmode\phantomsection\label{GELFAND-0158-U145804}\label{GELFAND-P001458:vanishing-cycles-degree}\textbf{\texttt{GELFAND-0158-U145804}} (vanishing-cycles-degree).\nopagebreak[4]\par
В определении после «т.е.» написано Φ с нижним f и малым тёмным верхним знаком, затем (𝓕)ₓ=0. Верхний знак наклонный, с утолщением и небольшим просветом. Область (1987,439,2130,535).\par
Кандидаты: \textbackslash{}Phi\_f\textasciicircum{}i(\textbackslash{}mathcal F)\_x=0; \textbackslash{}Phi\_f\textasciicircum{}\textbackslash{}bullet(\textbackslash{}mathcal F)\_x=0; \textbackslash{}Phi\_f\textasciicircum{}*(\textbackslash{}mathcal F)\_x=0. Чтение не установлено в дипломатическом исходнике.\par
Увеличение 10× не различает верхний индекс. Контекст комплекса исчезающих циклов допускает индекс или знак градуировки, но математическое ожидание не использовано для выбора. Независимый повторный просмотр также склоняется к i, но не устраняет неопределённость.\par
Повторная сверка: Повторно просмотрены исходник и увеличения 3×, 5× и 10×. Сохранены исходные наблюдения, варианты чтения и неуверенности; математическое рассуждение не заменяет буквальное чтение. Частичные воссоздания рисунков не утверждают неустановленных соединений или индексов.\par\medskip
\par\noindent\leavevmode\phantomsection\label{GELFAND-0158-U145805}\label{GELFAND-P001458:cotangent-cancellation}\textbf{\texttt{GELFAND-0158-U145805}} (cotangent-cancellation).\nopagebreak[4]\par
После df(X)⊂ и перед Ω¹ несколько толстых перекрещивающихся штрихов закрывают отменённый знак. Область (1912,945,2160,1100).\par
Кандидаты: не записаны. Чтение не установлено в дипломатическом исходнике.\par
Рассмотрено при 6×; не восстанавливаю зачёркнутый символ по ожидаемому обозначению кокасательного расслоения.\par
Повторная сверка: Повторно просмотрены исходник и увеличения 3×, 5× и 10×. Сохранены исходные наблюдения, варианты чтения и неуверенности; математическое рассуждение не заменяет буквальное чтение. Частичные воссоздания рисунков не утверждают неустановленных соединений или индексов.\par\medskip
\par\noindent\leavevmode\phantomsection\label{GELFAND-0158-U145806}\label{GELFAND-P001458:normalization-argument}\textbf{\texttt{GELFAND-0158-U145806}} (normalization-argument).\nopagebreak[4]\par
В строке «для морсовской ф-и f с df(x)=ν» далее отчётливо написано f(0)=0, с круглым нулём в аргументе. Область (275,1750,2180,2160), в особенности строка около y=2040–2120.\par
Кандидаты: f(0)=0 — буквальное чтение источника; f(x)=0 — возможная математическая нормировка, в источнике не утверждается. Чтение не установлено в дипломатическом исходнике.\par
Буквальное f(0)=0 сохранено после 3×. При этом предыдущая формула задаёт слой в x; возможно локальная координата x=0, возможно ошибка записи. Это математический вопрос, не неразборчивость цифры. Исправление в тело не внесено.\par
Повторная сверка: Повторно просмотрены исходник и увеличения 3×, 5× и 10×. Сохранены исходные наблюдения, варианты чтения и неуверенности; математическое рассуждение не заменяет буквальное чтение. Частичные воссоздания рисунков не утверждают неустановленных соединений или индексов.\par\medskip
\par\noindent\leavevmode\phantomsection\label{GELFAND-0158-U145807}\label{GELFAND-P001458:lower-detached-mark}\textbf{\texttt{GELFAND-0158-U145807}} (lower-detached-mark).\nopagebreak[4]\par
Над «если u∈𝓕\textasciicircum{}\{0,0\}» отдельно стоит небольшой наклонный знак с пересечением двух штрихов; напоминает 4 или знак вставки. Область (412,3090,800,3240).\par
Кандидаты: 4; знак вставки. Чтение не установлено в дипломатическом исходнике.\par
Рассмотрено при 6×. Не трактуется как номер страницы или доказанный номер пункта; сохранено расположение над началом строки.\par
Повторная сверка: Повторно просмотрены исходник и увеличения 3×, 5× и 10×. Сохранены исходные наблюдения, варианты чтения и неуверенности; математическое рассуждение не заменяет буквальное чтение. Частичные воссоздания рисунков не утверждают неустановленных соединений или индексов.\par\medskip
\subsection*{\texttt{GELFAND-P001459}}
\par\noindent\leavevmode\phantomsection\label{GELFAND-0158-U145901}\label{GELFAND-P001459:upper-cancellation}\textbf{\texttt{GELFAND-0158-U145901}} (upper-cancellation).\nopagebreak[4]\par
Перед «извращ. пучки» после «получаются» короткое слово многократно зачёркнуто. Область (1835,325,2030,392).\par
Кандидаты: не записаны. Чтение не установлено в дипломатическом исходнике.\par
Рассмотрено при 10×; исходные буквы слились с отменяющими штрихами. Само наличие отменённого слова сохранено.\par
Повторная сверка: Повторно просмотрены исходник и увеличения 3×, 5× и 10×. Сохранены исходные наблюдения, варианты чтения и неуверенности; математическое рассуждение не заменяет буквальное чтение. Частичные воссоздания рисунков не утверждают неустановленных соединений или индексов.\par\medskip
\par\noindent\leavevmode\phantomsection\label{GELFAND-0158-U145902}\label{GELFAND-P001459:bottom-cutoff}\textbf{\texttt{GELFAND-0158-U145902}} (bottom-cutoff).\nopagebreak[4]\par
Под длинной горизонтальной чертой у нижнего края короткого текстового фрагмента остаются небольшие верхние части штрихов; левый кластер напоминает М. Область (510,657,1520,712).\par
Кандидаты: М...; неразборчивые обрезанные штрихи. Чтение не установлено в дипломатическом исходнике.\par
Рассмотрено при 6×. Продолжение не выдумывается; все видимые обрезанные следы учтены. Остальная часть кадра ниже этого фрагмента пуста.\par
Повторная сверка: Повторно просмотрены исходник и увеличения 3×, 5× и 10×. Сохранены исходные наблюдения, варианты чтения и неуверенности; математическое рассуждение не заменяет буквальное чтение. Частичные воссоздания рисунков не утверждают неустановленных соединений или индексов.\par\medskip
\subsection*{\texttt{GELFAND-P001460}}
\par\noindent\leavevmode\phantomsection\label{GELFAND-0158-U146002}\label{GELFAND-P001460:stable-cancellation}\textbf{\texttt{GELFAND-0158-U146002}} (stable-cancellation).\nopagebreak[4]\par
В строке «Стаб. пучки и» следующее слово зачёркнуто несколькими толстыми горизонтальными штрихами; уверенных целых букв не сохранилось. Область (870,2890,1080,2960).\par
Кандидаты: не записаны. Чтение не установлено в дипломатическом исходнике.\par
Рассмотрено при 10×, также строка (445,2870,1090,2980) при 3×. Чёткой закрывающей кавычки после зачёркивания нет; она не добавлена.\par
Повторная сверка: Повторно просмотрены исходник и увеличения 3×, 5× и 10×. Сохранены исходные наблюдения, варианты чтения и неуверенности; математическое рассуждение не заменяет буквальное чтение. Частичные воссоздания рисунков не утверждают неустановленных соединений или индексов.\par\medskip
\par\noindent\leavevmode\phantomsection\label{GELFAND-0158-U146003}\label{GELFAND-P001460:quantum-title-mark}\textbf{\texttt{GELFAND-0158-U146003}} (quantum-title-mark).\nopagebreak[4]\par
Между «Квант. группы» и «(по Дринфельду)» стоит малый плотный знак над строкой: верхние горизонтальные штрихи и два нижних выступа. Область (830,2985,928,3055).\par
Кандидаты: зачёркнутые закрывающие кавычки; звёздочка. Чтение не установлено в дипломатическом исходнике.\par
Рассмотрено при 10×. Знак похож на закрашенные кавычки, но не заменён уверенной кавычкой или знаком сноски. После скобок в источнике имеются отдельные закрывающие кавычки.\par
Повторная сверка: Повторно просмотрены исходник и увеличения 3×, 5× и 10×. Сохранены исходные наблюдения, варианты чтения и неуверенности; математическое рассуждение не заменяет буквальное чтение. Частичные воссоздания рисунков не утверждают неустановленных соединений или индексов.\par\medskip
\subsection*{\texttt{GELFAND-P001461}}
\par\noindent\leavevmode\phantomsection\label{GELFAND-0158-U146101}\label{GELFAND-P001461:surface-sketch}\textbf{\texttt{GELFAND-0158-U146101}} (surface-sketch).\nopagebreak[4]\par
Рисунок слева в средней части страницы, исходные координаты (170,1735,525,1985): внутренний овальный контур с двумя маленькими петлеобразными штрихами, наружный широкий контур с пунктирными участками. Подпись «рим. пов. ⊂ сфера» читается и отдельно передана в тексте.\par
Кандидаты: Две внутренние петлеобразные отметки, соединение верхней с контуром не установлено; Условное изображение ручек поверхности. Чтение не установлено в дипломатическом исходнике.\par
Осмотрен оригинал 2408×3480 и увеличенный в 3 раза crop (170,1240,1310,2075). Точную топологию, соединение верхней петли и неодинаковые типы линий невозможно уверенно восстановить; SourceFigure сохраняет привязку к оригиналу без придуманной топологии.\par
Повторная сверка: Повторно просмотрены исходник и увеличения 3×, 5× и 10×. Сохранены исходные наблюдения, варианты чтения и неуверенности; математическое рассуждение не заменяет буквальное чтение. Частичные воссоздания рисунков не утверждают неустановленных соединений или индексов.\par\medskip
\subsection*{\texttt{GELFAND-P001462}}
\par\noindent\leavevmode\phantomsection\label{GELFAND-0158-U146201}\label{GELFAND-P001462:antipode-cancelled-sign}\textbf{\texttt{GELFAND-0158-U146201}} (antipode-cancelled-sign).\nopagebreak[4]\par
После номера «3.» перед «S:A→A» имеется густо перечёркнутый одиночный знак. Исходные координаты (345,1260,430,1388).\par
Кандидаты: Δ; и. Чтение не установлено в дипломатическом исходнике.\par
Оригинал и crop 10× просмотрены. Несколько диагональных штрихов перекрывают знак; ни один кандидат не утверждается. Зачёркнутый участок не пропущен. Независимый повторный просмотр root при 10× также не позволил различить знак однозначно.\par
Повторная сверка: Повторно просмотрены исходник и увеличения 3×, 5× и 10×. Сохранены исходные наблюдения, варианты чтения и неуверенности; математическое рассуждение не заменяет буквальное чтение. Частичные воссоздания рисунков не утверждают неустановленных соединений или индексов.\par\medskip
\par\noindent\leavevmode\phantomsection\label{GELFAND-0158-U146202}\label{GELFAND-P001462:antipode-formula}\textbf{\texttt{GELFAND-0158-U146202}} (antipode-formula).\nopagebreak[4]\par
Формула (4) явно записана S(X\textasciicircum{}±)=q\textasciicircum{}\{H/2\}X\textasciicircum{}±q\textasciicircum{}\{-H/2\}, без начального минуса в правой части. Исходный crop (410,1390,1290,1610), увеличение 5×.\par
Кандидаты: S(X\textasciicircum{}±)=q\textasciicircum{}\{H/2\}X\textasciicircum{}±q\textasciicircum{}\{-H/2\} — буквальная запись источника. Чтение не установлено в дипломатическом исходнике.\par
Математическая сверка с показанной ниже диаграммой антипода и с S(H)=−H выявляет возможную ошибку знака/нормировки в источнике. Это не сомнение в наличии видимого минуса: его перед q нет. Запись сохранена, исправление не внесено.\par
Повторная сверка: Повторно просмотрены исходник и увеличения 3×, 5× и 10×. Сохранены исходные наблюдения, варианты чтения и неуверенности; математическое рассуждение не заменяет буквальное чтение. Частичные воссоздания рисунков не утверждают неустановленных соединений или индексов.\par\medskip
\par\noindent\leavevmode\phantomsection\label{GELFAND-0158-U146203}\label{GELFAND-P001462:cancelled-generators}\textbf{\texttt{GELFAND-0158-U146203}} (cancelled-generators).\nopagebreak[4]\par
После «образующих» густо зачёркнут фрагмент, в котором различимы формы, похожие на K, X\textasciicircum{}±; правее остаются перекрытые петлеобразные штрихи. Исходные координаты (1530,2410,1915,2580), review crop (1340,2370,1925,2610), 5×.\par
Кандидаты: K,X\textasciicircum{}±, затем неразборчивое; K,X\textasciicircum{}±, q…. Чтение не установлено в дипломатическом исходнике.\par
Более поздняя строка «K=q\textasciicircum{}\{H/2\}, K,X\textasciicircum{}±» и нижнее зачёркивание K,X\textasciicircum{}± читаются отдельно и не использованы для восстановления этого зачёркнутого фрагмента.\par
Повторная сверка: Повторно просмотрены исходник и увеличения 3×, 5× и 10×. Сохранены исходные наблюдения, варианты чтения и неуверенности; математическое рассуждение не заменяет буквальное чтение. Частичные воссоздания рисунков не утверждают неустановленных соединений или индексов.\par\medskip
\par\noindent\leavevmode\phantomsection\label{GELFAND-0158-U146204}\label{GELFAND-P001462:commutator-normalization}\textbf{\texttt{GELFAND-0158-U146204}} (commutator-normalization).\nopagebreak[4]\par
Внизу записано X\textasciicircum{}+X\textasciicircum{}−−X\textasciicircum{}−X\textasciicircum{}+=(K−K\textasciicircum{}\{−1\})/(q−q\textasciicircum{}\{−1\}), выше K=q\textasciicircum{}\{H/2\}. Предыдущая страница GELFAND-P001461 заканчивается [X\textasciicircum{}+,X\textasciicircum{}−]=sh((h/2)H)/sh(h/2), q=e\textasciicircum{}h. Исходный crop нижней формулы (370,2925,1310,3180), 5×.\par
Кандидаты: (K−K\textasciicircum{}\{−1\})/(q−q\textasciicircum{}\{−1\}) — буквальная запись источника. Чтение не установлено в дипломатическом исходнике.\par
При подстановке предыдущих обозначений знаменатели двух видимых формул различаются: sh(h/2) даёт q\textasciicircum{}\{1/2\}−q\textasciicircum{}\{−1/2\}. Изменение соглашения в пределах этих строк не обозначено. Источник воспроизведён без исправления знаменателя.\par
Повторная сверка: Повторно просмотрены исходник и увеличения 3×, 5× и 10×. Сохранены исходные наблюдения, варианты чтения и неуверенности; математическое рассуждение не заменяет буквальное чтение. Частичные воссоздания рисунков не утверждают неустановленных соединений или индексов.\par\medskip
\subsection*{\texttt{GELFAND-P001463}}
\par\noindent\leavevmode\phantomsection\label{GELFAND-0158-U146301}\label{GELFAND-P001463:r-coefficient}\textbf{\texttt{GELFAND-0158-U146301}} (r-coefficient).\nopagebreak[4]\par
Формула (8), числитель коэффициента: q\textasciicircum{}n, знаменатель [n]\_q!, затем e\textasciicircum{}n⊗f\textasciicircum{}n. Координаты исходного PNG: x=1040–1230, y=1160–1400; просмотр 3× и 6×.\par
Кандидаты: q\textasciicircum{}n/[n]\_q!. Чтение не установлено в дипломатическом исходнике.\par
Степень n читается прямо. Дополнительных степеней или множителей в этой записи не видно; они не добавлены. Математическая полнота рукописного коэффициента оставлена для проверки.\par
Повторная сверка: Повторно просмотрены исходник и увеличения 3×, 5× и 10×. Сохранены исходные наблюдения, варианты чтения и неуверенности; математическое рассуждение не заменяет буквальное чтение. Частичные воссоздания рисунков не утверждают неустановленных соединений или индексов.\par\medskip
\par\noindent\leavevmode\phantomsection\label{GELFAND-0158-U146302}\label{GELFAND-P001463:cancelled-before-membership}\textbf{\texttt{GELFAND-0158-U146302}} (cancelled-before-membership).\nopagebreak[4]\par
Перед скобкой (∈ U\_q(SU(2))\textasciicircum{}\{⊗3\}) справа от первой строки (10) несколько густо зачёркнутых знаков. Исходные координаты: x=1410–1525, y=1900–2045; просмотр 3× и 10×.\par
Кандидаты: не записаны. Чтение не установлено в дипломатическом исходнике.\par
Перечёркивающие штрихи не дают надёжно отделить исходные буквы; зачёркивание отмечено в тексте, содержание не восстановлено.\par
Повторная сверка: Повторно просмотрены исходник и увеличения 3×, 5× и 10×. Сохранены исходные наблюдения, варианты чтения и неуверенности; математическое рассуждение не заменяет буквальное чтение. Частичные воссоздания рисунков не утверждают неустановленных соединений или индексов.\par\medskip
\par\noindent\leavevmode\phantomsection\label{GELFAND-0158-U146303}\label{GELFAND-P001463:exercise-harmonics}\textbf{\texttt{GELFAND-0158-U146303}} (exercise-harmonics).\nopagebreak[4]\par
В приписке после «по сферич.» стоит сокращение «гар-…м»: между дефисом и последним м виден небольшой петлеобразный штрих. Исходные координаты: x=1890–2175, y=1030–1105; просмотр 3× и независимый просмотр 8×.\par
Кандидаты: гар-ям; гар-м. Чтение не установлено в дипломатическом исходнике.\par
Основная строка и помета «начать с SU(2)» прочитаны; спорным остаётся только сокращение. Сокращение не развёрнуто.\par
Повторная сверка: Повторно просмотрены исходник и увеличения 3×, 5× и 10×. Сохранены исходные наблюдения, варианты чтения и неуверенности; математическое рассуждение не заменяет буквальное чтение. Частичные воссоздания рисунков не утверждают неустановленных соединений или индексов.\par\medskip
\subsection*{\texttt{GELFAND-P001464}}
\par\noindent\leavevmode\phantomsection\label{GELFAND-0158-U146401}\label{GELFAND-P001464:definition-separator}\textbf{\texttt{GELFAND-0158-U146401}} (definition-separator).\nopagebreak[4]\par
В первой строке между R\textasciicircum{}\{j₁j₂\} и V\textasciicircum{}\{j₁\} два утолщённых элемента, один над другим, с дополнительными горизонтальными штрихами. Исходные координаты: x=490–567, y=142–235; просмотр 3× и 10×.\par
Кандидаты: двоеточие; двоеточие поверх исправленного знака. Чтение не установлено в дипломатическом исходнике.\par
Формула дана с кандидатом «:» и явной пометой. Первоначальный исправленный знак не установлен.\par
Повторная сверка: Повторно просмотрены исходник и увеличения 3×, 5× и 10×. Сохранены исходные наблюдения, варианты чтения и неуверенности; математическое рассуждение не заменяет буквальное чтение. Частичные воссоздания рисунков не утверждают неустановленных соединений или индексов.\par\medskip
\par\noindent\leavevmode\phantomsection\label{GELFAND-0158-U146402}\label{GELFAND-P001464:cancelled-coefficient}\textbf{\texttt{GELFAND-0158-U146402}} (cancelled-coefficient).\nopagebreak[4]\par
После «q-аналоги» перед «коэфф. Клебша-Гордана» один или несколько густо зачёркнутых буквенных штрихов. Исходные координаты: x=677–746, y=1110–1205; просмотр 3× и 10×.\par
Кандидаты: не записаны. Чтение не установлено в дипломатическом исходнике.\par
Надёжных кандидатов для исходного зачёркнутого знака не получено; он не заменён ожидаемым словом.\par
Повторная сверка: Повторно просмотрены исходник и увеличения 3×, 5× и 10×. Сохранены исходные наблюдения, варианты чтения и неуверенности; математическое рассуждение не заменяет буквальное чтение. Частичные воссоздания рисунков не утверждают неустановленных соединений или индексов.\par\medskip
\par\noindent\leavevmode\phantomsection\label{GELFAND-0158-U146403}\label{GELFAND-P001464:sum-mark}\textbf{\texttt{GELFAND-0158-U146403}} (sum-mark).\nopagebreak[4]\par
В (15) над суммой справа маленький кружок, с утолщёнными вертикальной и горизонтальной чертами и диагональным штрихом в нижней левой части. Исходные координаты: x=1015–1085, y=1245–1310; просмотр 3× и 10×.\par
Кандидаты: ⊕; ⊗, исправленное на ⊕. Чтение не установлено в дипломатическом исходнике.\par
В формуле сохранён читаемый финальный плюс в кружке; возможное исправление первоначального знака остаётся неразрешённым. Знак суммы не заменён единым символом прямой суммы.\par
Повторная сверка: Повторно просмотрены исходник и увеличения 3×, 5× и 10×. Сохранены исходные наблюдения, варианты чтения и неуверенности; математическое рассуждение не заменяет буквальное чтение. Частичные воссоздания рисунков не утверждают неустановленных соединений или индексов.\par\medskip
\par\noindent\leavevmode\phantomsection\label{GELFAND-0158-U146404}\label{GELFAND-P001464:coefficient-range}\textbf{\texttt{GELFAND-0158-U146404}} (coefficient-range).\nopagebreak[4]\par
Внизу слева написано A с верхними индексами m′₁…m′\_N и нижними m₁…m\_L. Исходные координаты: x=370–813, y=2860–3090; просмотр 3× и 5×.\par
Кандидаты: A\textasciicircum{}\{m′₁…m′\_N\}\_\{m₁…m\_L\}. Чтение не установлено в дипломатическом исходнике.\par
В строке отображения A выше N стоит у исходных множителей, L — у конечных; диаграммы также имеют нижние j\_N и верхние j′\_L. В записи коэффициента конечные индексы N/L стоят в обратных рядах. Видимая запись сохранена; согласование обозначений не выполнено.\par
Повторная сверка: Повторно просмотрены исходник и увеличения 3×, 5× и 10×. Сохранены исходные наблюдения, варианты чтения и неуверенности; математическое рассуждение не заменяет буквальное чтение. Частичные воссоздания рисунков не утверждают неустановленных соединений или индексов.\par\medskip
\par\noindent\leavevmode\phantomsection\label{GELFAND-0158-U146405}\label{GELFAND-P001464:diagram-02}\textbf{\texttt{GELFAND-0158-U146405}} (diagram-02).\nopagebreak[4]\par
Нижняя правая схема: прямоугольник A, по три вертикальных вывода сверху и снизу без стрелок. Вверху j′₁,j′₂,j′\_L, над ними m′₁,m′₂,…,m′ с густо исправленным последним индексом; внизу j₁,j₂,j\_N и m₁,m₂,m\_N. Схема x=1370–2035,y=2615–3445; спорный индекс x=1870–2010,y=2640–2765. Просмотр схемы 3×, индекса 10×.\par
Кандидаты: верхний правый индекс L поверх N; верхний правый индекс L поверх неразборчивого знака; верхний правый индекс N с правкой. Чтение не установлено в дипломатическом исходнике.\par
Неясна исходная буква под густой правкой у верхней правой подписи. Поэтому вся эта схема оставлена SourceFigure с описанием всех видимых выводов и подписей. Первая схема (17) с однозначными подписями передана inline TikZ.\par
Повторная сверка: Повторно просмотрены исходник и увеличения 3×, 5× и 10×. Сохранены исходные наблюдения, варианты чтения и неуверенности; математическое рассуждение не заменяет буквальное чтение. Частичные воссоздания рисунков не утверждают неустановленных соединений или индексов.\par\medskip
